\vfill\pagebreak
\section{Character and String types}
\label{sec:types:chars_strings}

A character type holds one textual symbol -- a letter, a digit, a punctuation
mark. Put characters in sequence and the result is a string, which is how a
program holds anything longer than a symbol: a word, a line, a file.

\subsection{Character literal}
\label{sec:types:char_literal}

A character literal is a symbol between single quotes, \token{'}. Three types
hold one: \token{c8}, \token{c16} and \token{c32}, for UTF-8, UTF-16 and UTF-32
respectively. The encoding is what decides how many bits a symbol takes, and once
strings enter the picture that choice stops being a detail --
Section~\ref{sec:types:encoding} is about little else.

\begin{lstlisting}[style=coloredverbatim, label=lst:types:char_literal, caption=Example of character literals, escapechar=@]
let x = 'A';
let y = '@\ensuremath{\color{black!30!applegreen} \texttt{\Omega{}}}@'c16;
let z = '@\emoji{slightly-smiling-face}@'c32;
\end{lstlisting}


Listing~\ref{lst:types:char_literal} stores three of them. A suffix names the
type, as it does for numbers: \token{'Ω'c16} is a \token{c16} and
\token{'🙂'c32} a \token{c32}. An unsuffixed literal such as \token{'A'} is
always a \token{c8}, with no exception made for wide characters.

Which is why \token{'🙂'} on its own does not compile. That symbol needs a full
32 bits, and an unsuffixed literal offers it 8.

\begin{lstlisting}[style=coloredverbatimError, label=lst:types:char_encode_error, escapechar=@]
let x = '@\hb{\emoji{slightly-smiling-face}}@'; // Error : malformed literal, number of c8 is 4 
\end{lstlisting}%

Some characters cannot be typed inside a literal: a line break would end the
line, and a quote would end the literal. These are written as escape sequences
instead, always introduced by a backslash. Ymir accepts the ones below, and all
of them except \tokennolst{\textbackslash{}u\{code\}} -- whose width depends on
the code point it names -- fit in a single UTF-8 character.

\begin{center}
  \begin{tabular}{ll}
    Value & Content\\[0pt]
    \hline
    \tokennolst{\textbackslash{}a} & Alert beep, (Bell)\\[0pt]
    \tokennolst{\textbackslash{}b} & Backspace\\[0pt]
    \tokennolst{\textbackslash{}f} & Page break\\[0pt]
    \tokennolst{\textbackslash{}n} & New line\\[0pt]
    \tokennolst{\textbackslash{}r} & Carriage return\\[0pt]
    \tokennolst{\textbackslash{}t} & Horizontal tab\\[0pt]
    \tokennolst{\textbackslash{}v} & Vertical tab\\[0pt]
    \tokennolst{\textbackslash{}\textbackslash{}} & Backslash\\[0pt]
    \tokennolst{\textbackslash{}'} & Apostrophe\\[0pt]
    \tokennolst{\textbackslash{}"} & Double quotation mark\\[0pt]
    \tokennolst{\textbackslash{}u\{code\}} & Unicode\\[0pt]
  \end{tabular}
  \captionof{table}{\label{tab:types:escape_chars} Escape characters}
\end{center}

The last of them, \tokennolst{\textbackslash{}u\{code\}}, is the way to write a
character by its Unicode code point rather than by the symbol itself -- useful
when the symbol is hard to type, or hard to tell apart from another. The code
point of \token{'🙂'} is \token{0x1F642}, so it may equally be written
\token{'\u\{0x1F642\}'c32}.

\subsection{String type}

A string holds text: characters laid out in memory one after another, in some
encoding. A single character type is fixed in width, but a string is not -- what
it costs in memory depends on which symbols are in it.

\subsection{String literal}

A string literal is a run of symbols between double quotes, \token{"}. Its type
is written in brackets, and reads as an array of characters: \token{\[c8\]} is a
UTF-8 string of any length, \token{\[c16 ; 5\]} a UTF-16 string of exactly five.
The notation states the encoding and the length constraint in one place. Arrays
and slices are the subject of a later chapter; all that matters here is that both
store their elements consecutively in memory.

String literals take suffixes too -- \token{s8}, \token{s16} and \token{s32} --
and default to UTF-8 without one.

\begin{lstlisting}[style=coloredverbatim, label=lst:types:string_literal, caption=Example of string literals, escapechar=@]
let greetings = "Hello World!";
let utf8_unicode = "@\korean{{\color{black!30!applegreen} 안녕하세요 세계}}@";
let utf32_unicode = "@\japan{{\color{black!30!applegreen} こんにちは 世界}}@"s32;
\end{lstlisting}

\subsection{A small word on encoding}
\label{sec:types:encoding}

A character may take anything from one to four bytes, as
Section~\ref{sec:types:char_literal} noted. The characters of the ASCII table
take one, which is why they sit in a \token{c8}, or in a \token{[c8]} string,
with nothing clever going on. The \token{greetings} variable of
Listing~\ref{lst:types:string_literal} is one of these: plain ASCII, one byte per
character, laid out in order, starting with \token{'H'} -- which is $72$.

\input{chapters/chapter2/figures/hello_string_figure}


Everything outside ASCII goes through Unicode, which assigns a code point to
every character in every writing system. The three encodings differ in how they
store those code points, and the trade-off is always the same: simplicity against
space.

UTF-32 takes the simple end. Every character gets four bytes, whether it needs
them or not, and the surplus is zeros. Nothing is ever ambiguous, and a good deal
of memory is wasted.



The following Japanese text \colorbox{background!90!gray!60}{\japan{\color{black!30!applegreen} "こんにちは"}}
is composed of five characters, each represented as a UTF‑32
Unicode code point: \token{0x3053 0x3093 0x306b 0x3061 0x306f}. The first
character, \colorbox{background!90!gray!60}{\japan{\color{black!30!applegreen} こ}}, has the Unicode code point
\token{0x3053}, which requires only 14 bits. This means it could be stored
within a \token{c16} type. However, under UTF‑32 encoding, every character is
allocated a full 4 bytes, leaving two bytes filled with zeros. The same
applies to each of the five characters in the word, resulting in a total of 20
bytes used, even though only 10 bytes contain meaningful information. On the
other hand, the character \token{🙂} has the Unicode code point
\token{0x1f642}, which requires 17 bits and therefore no longer fits inside a
\token{c16}. As a result, UTF‑16 and UTF‑32 encodings of that character take
the same number of bytes to store it in memory.



UTF‑8 encodes Unicode characters by using a variable number of bytes for each
symbol. The first byte contains information that indicates how many subsequent
bytes belong to the same character. For example, in the mixed English and Korean
text \colorbox{background!90!gray!60}{\color{black!30!applegreen}{"\texttt{Hello is }\korean{안녕하세요}"}}, the segment \colorbox{background!90!gray!60}{\color{black!30!applegreen}{"\texttt{Hello is }"}} uses one
byte per character, while the segment \colorbox{background!90!gray!60}{\korean{\color{black!30!applegreen}{"안녕하세요"}}} requires three bytes
for each character. The UTF‑8 encoding of the entire string therefore occupies
24 bytes of memory, compared to 56 bytes in UTF‑32 and 28 bytes in UTF‑16.


The saving has a catch, and it is the important part of this section: the length
of a UTF-8 string is a count of bytes, not of characters.
\colorbox{background!90!gray!60}{\korean{\color{black!30!applegreen}{"안녕하세요"}}} is
five characters and fifteen bytes, and it is fifteen that a \token{\[c8\]}
reports. UTF-16 has the same problem in milder form, its characters taking two
bytes or four depending on the code point.

\begin{lstlisting}[style=coloredverbatim, label=lst:types:encoding_example, caption=Encoding consequence on array length, escapechar=@]
let utf8_str = "@\korean{\color{black!30!applegreen} 안녕하세요 세계}@";
let utf16_str = "@\korean{\color{black!30!applegreen} 안녕하세요 세계}@"s16;
let utf32_str = "@\korean{\color{black!30!applegreen} 안녕하세요 세계}@"s32;

println(utf8_str.len); // 22, or 22 bytes
println(utf16_str.len); // 8, or 8 * 2 => 16 bytes
println(utf32_str.len); // 8, or 8 * 4 => 32 bytes
\end{lstlisting}

\notebox{ For all these encoding reasons, it becomes clear why the character
  \emoji{slightly-smiling-face} cannot be represented using a single \token{c8}. In fact, the
  length of \token{"🙂"} is reported as 4, since UTF‑8 requires four bytes to
  encode this symbol. Encoding is not a trivial topic. It will be explored in
  greater detail in a later chapter. What is important to understand here is
  that the apparent length of a string does not always correspond to the number
  of characters it contains. When in doubt, converting the string into UTF‑32
  will reveal the true character count, as UTF‑32 consistently encodes each
  character with a fixed length of four bytes. }

\begin{lstlisting}[style=coloredverbatim, label=lst:types:converting_encoding, caption=String encoding conversion, escapechar=@]
use std::conv; // to import 'to'

let utf8_str = "@\korean{\color{black!30!applegreen} 안녕하세요 세계}@";

let utf32_str = to!{[c32]}(utf8_str); // converting UTF8 into UTF32

assert(utf8_str.len == 22);
assert(utf32_str.len == 8);
\end{lstlisting}
